\chapter{Orthostatic Hypotension}

\section*{Definition}
A sustained decrease in systolic blood pressure of at least 20~mm Hg or diastolic blood pressure of 10~mm Hg within 3~minutes of standing or upright tilt, often accompanied by cerebral hypoperfusion symptoms.

\section*{What Happens in Your Body}
Upon standing, gravity causes pooling of 500--800~mL of blood in the lower extremities and splanchnic circulation. Inadequate compensatory narrowing of blood vessels or chronotropic response, due to autonomic dysfunction, volume depletion, or medication effect, leads to a transient drop in cerebral perfusion pressure and clinical symptoms.

\section*{Causes}
\begin{itemize}
  \item \textbf{Volume depletion:} Dehydration, hemorrhage, diuretics, burns, adrenal insufficiency
  \item \textbf{Neurogenic:} Parkinson disease, diabetes mellitus, amyloidosis, pure autonomic failure, multiple system atrophy
  \item \textbf{Cardiovascular:} Aortic stenosis, cardiomyopathy, arrhythmias, myocardial infarction
  \item \textbf{Medications:} Antihypertensives, alpha-blockers, tricyclic antidepressants, antipsychotics, vasodilators
  \item \textbf{Miscellaneous:} Prolonged bed rest, aging, anemia, alcohol use, postprandial splanchnic pooling
\end{itemize}

\section*{When to See a Doctor}
See your doctor if:
\begin{itemize}
  \item Recurrent syncope or near-syncope upon standing
  \item Falls, especially in older adults
  \item Accompanying chest pain, palpitations, or severe headache
\end{itemize}

\section*{Related Symptoms}
\begin{itemize}
  \item Dizziness, lightheadedness, presyncope, syncope
  \item Blurred vision, weakness, fatigue, cognitive slowing
  \item Neck or shoulder pain (coat-hanger distribution)
\end{itemize}
\textbf{Important:} Symptoms are often worse in the morning, after meals, with heat exposure, or with physical exertion.

\section*{Self-Care / Home Management}
\begin{itemize}
  \item Increase fluid and salt intake if not contraindicated by heart failure or hypertension
  \item Rise slowly from lying to sitting to standing; dangle legs before standing
  \item Wear compression stockings (30--40~mm Hg) and abdominal binders
  \item Elevate head of bed 10--30\textdegree{} to reduce nocturnal diuresis
  \item Avoid large carbohydrate-heavy meals, alcohol, and hot showers
\end{itemize}

\section*{How Your Doctor Will Evaluate You}
\begin{itemize}
  \item Orthostatic vital signs: supine and standing BP/HR at 1 and 3~minutes
  \item Basic metabolic panel, CBC, and thyroid function to exclude secondary causes
  \item 24-hour ambulatory BP monitoring and tilt-table testing if suspicion is high
  \item Autonomic reflex testing and plasma norepinephrine levels for neurogenic cases
\end{itemize}

\section*{Treatment Options}
\begin{itemize}
  \item Correct volume depletion and discontinue offending medications if possible
  \item Pharmacologic: midodrine (alpha-agonist), fludrocortisone (mineralocorticoid), droxidopa for neurogenic OH
  \item Non-pharmacologic: counterpressure maneuvers (leg crossing, squatting, muscle pumping)
  \item Lifestyle: salt tablets (1--2~g TID), adequate hydration (2--3~L daily), compression garments
\end{itemize}

\clearpage
